% Physical pages 3 and 4 of the concise study: the thematic orientation, in the
% order of the Mass, with what each text says, keeps and omits. Every claim and
% quotation comes from the reviewed expansive study. The developed comparison
% of the four interpretations begins on the fifth page.

\runninghead{Themes and Movement}
\section{The Propers: Themes and Movement}
\zlabel{triptych:concise:themes:start}

Three texts of this Mass say that something is near. St Paul tells the Romans
that ``our salvation is nearer than when we believed'' (Rom 13:11); in the
Gospel Christ tells his disciples that their redemption ``is at hand,'' and by
the parable of the fig tree that the kingdom of God is near (Lk 21:28, 31); and the
Postcommunion looks to the solemnities ``to come'' of our restoration. Between
the lessons the Alleluia asks God to ``grant us thy salvation.'' The four
interpretations of the Mass differ in what they find near. In the first,
St Gregory the Great reads the Gospel's signs as real blows on an aging
world, and the lifted head as a mind gladdened because a world it does not
love is ending; St Cyril of Alexandria says that the redemption is the body's
resurrection and creation made new. In the second, St Ambrose and St
Augustine read the failing lights of heaven as figures of the Church under
persecution, and find Christ's coming already under way: daily in his body
the Church, for Augustine, and, Ambrose suggests with a ``perhaps,'' in each
one who receives him. In the third, St John Chrysostom reads Paul's hour as
the hour of the end, nearer every day, and St Thomas Aquinas finds salvation
near in time, by the growth of charity, and, as the Church reads Paul in
Advent, in the first coming. In the fourth, Augustine, preaching on the
Alleluia's psalm, says that God gave Christ as man and will give the same
Christ as God, and St Robert Bellarmine reads the psalm as a prayer for the
Saviour's coming.

The Mass stands first in the Missal, headed by its Roman station, St Mary
Major. Blessed Ildefonso Schuster recalls that this basilica has venerated,
since Pope Sixtus III in the fifth century, a reproduction of the crib of
Bethlehem, and he takes the station to point the faithful toward the manger
at the start of the season (\work{The Sacramentary}, vol.~1, p.~321). The
Introit opens with \latin{Ad te levávi ánimam meam}, ``To thee have I lifted
up my soul.'' Psalm 24 is a prayer for
guidance and forgiveness. After the verses the Introit sings, the psalmist
confesses ``the sins of my youth and my ignorances'' (v.~7), and he ends by
praying for the whole people, ``Deliver Israel, O God, from all his
tribulations'' (v.~22). The man who lifts his soul is in danger, burdened by
his sins and mocked by enemies. The antiphon keeps three of his acts: he lifts
his soul, he trusts, and he appeals to a rule he has learned, that none who
wait for God are put to shame. Where the Vulgate reads \latin{qui sustinent
te}, the chant reads \latin{qui te exspéctant}, ``who wait for thee,'' as
Augustine's psalter did. Augustine hears in the psalm Christ speaking in the
person of the Church (\work{En.\ in Ps.}~24.1--2). The same antiphon returns
at the offering and the Gradual takes its words from the same psalm, so that
Psalm 24 is sung three times; Schuster calls it ``the typical psalm of the
first Sunday of Advent'' (p.~322).

The Collect ends \latin{Qui vivis et regnas cum Deo Patre}, ``who livest and
reignest with God the Father,'' and so is addressed to the Son. It asks him to
stir up his power,
\latin{Excita \dots\ poténtiam tuam}, and then simply \latin{et veni}, ``and
come.'' It names what it fears, \latin{imminéntibus peccatórum nostrórum
perículis}, ``the threatening dangers of our sins'': the dangers hang over
us, and our sins have made them. Two verbs follow, each with Christ's action
beside it: \latin{te protegénte éripi}, to be snatched away with him
protecting, and \latin{te liberánte salvári}, to be saved with him setting
free.

The Epistle comes from the moral part of Romans. Just before it Paul has
said that ``love \dots\ is the fulfilling of the law'' (13:10), and the
Vulgate goes on, ``And that, knowing the season'': love is to be practised by
people who know what time it is. The Missal opens instead \latin{Fratres:
Sciéntes}, ``Brethren, knowing,'' and the lesson stands alone. Paul's argument
has three steps. It is the hour to rise from sleep, because salvation is
nearer than when we first believed; the night is past and the day at hand, so
the works of darkness are to be thrown off and the armour of light put on; and
to walk becomingly as in the day means leaving revelling and drunkenness, the
bed and wantonness, strife and envy. The verbs are of clothing:
\latin{abiciámus}, let us cast off, \latin{induámur}, let us put on,
\latin{induímini}, put ye on. The lesson stops in the middle of verse~14,
before ``and make not provision for the flesh in its concupiscences,'' so
that its last words are ``put ye on the Lord Jesus Christ.'' Augustine read
these verses in a garden, after a child's voice chanted ``Take up and read,''
and was converted (\work{Confessions} VIII.12.29).

The Gradual returns to Psalm 24 for the rule the Introit cited and adds a
final \latin{Dómine}, so that the rule becomes an address: ``None who wait for
thee shall be confounded, O Lord.'' Bellarmine explains why the psalmist can
state it: ``we have learned by long experience, from the examples of our
ancestors, and from your own promises,'' that those who patiently expect God's
help ``were never disappointed'' (\work{Explanatio}, Ps 24, on vv.~2--3). The
verse asks God to make his ways known, \latin{notas fac mihi}, where the
Introit's verse has \latin{demónstra}, ``shew.''

Psalm 84, from which the Alleluia comes, begins with thanks for a
restoration already given, ``thou hast turned away the captivity of Jacob,''
and then asks for it again: ``Convert us, O God our saviour: and turn off thy
anger from us'' (vv.~2, 5). The Alleluia takes the verse that asks, ``Shew us,
O Lord, thy mercy; and grant us thy salvation'' (v.~8). Its word
for salvation, \latin{salutáre}, is an adjective used as a noun, ``thy saving
thing,'' and Augustine reads it as a person: \latin{salutare tuum, Christum
tuum}, ``thy salvation, thy Christ'' (\work{En.\ in Ps.}~84.9).

The Gospel is the end of Christ's last discourse in St Luke, spoken in the
Temple in the days before his Passion. It began when some admired the
Temple's stones; just before the lesson Christ foretells Jerusalem compassed
with armies and ``trodden down by the Gentiles'' (21:20--24). The Missal
begins at the signs with ``At that time,'' and Augustine and Cyril of
Alexandria both take Luke to pass at this point from the fall of the city to
the end (Letter 199.27; \work{Comm.\ in Luc.}, sermon 139). The lesson has three parts: the signs in the heavens, with the
nations in distress and men withering for fear; the Son of man coming
\latin{in nube}, ``in a cloud,'' with power and majesty; and the command
\latin{respícite, et leváte cápita vestra}, ``look up and lift up your
heads,'' because redemption is near. The Introit's \latin{levávi} and the
Gospel's \latin{leváte} are the same verb. A parable follows: when the fig
tree and all the trees put out their fruit, summer is near, and so is the
kingdom of God when these things happen. Two sayings close the lesson. The
first is introduced by \latin{Amen}: this generation will not pass till all
is done; then heaven and earth will pass, but Christ's words will not. Just
after it Christ warns
against hearts ``overcharged with surfeiting and drunkenness'' and commands,
``Watch ye, therefore, praying at all times'' (vv.~34, 36).

At the offering the choir sings the Introit's antiphon again, without its
verse. The Secret's subject is \latin{Hæc sacra}, ``these holy things,'' the
gifts on the altar. It asks that, having cleansed us \latin{poténti
virtúte}, ``by their mighty power,'' they make us come \latin{purióres},
``purer,'' \latin{ad suum princípium}, to their own source. Schuster finds in
it the idea that runs through the Roman epiklesis: the sacrifice's virtue is
asked to purify our minds ``that we may worthily receive him who is its
Author'' (p.~322).

The Communion returns to Psalm 84 for the verse that closes its promise: ``The
Lord will give goodness: and our earth shall yield her fruit'' (v.~13). In the
psalm it follows ``Truth is sprung out of the earth: and justice hath looked
down from heaven'' (v.~12), and ``Justice shall walk before him'' follows it
(v.~14). The chant leaves out the Vulgate's opening
\latin{Etenim}, ``For,'' and keeps both halves of the verse, heaven's gift and
the earth's yield. Augustine's psalter read \latin{suavitatem},
``sweetness,'' where the Missal has \latin{benignitátem}, ``goodness.'' Who or
what the fruitful earth is divides the commentators.

The Postcommunion borrows Ps 47:10 (48:10), ``We have received thy mercy, O
God, in the midst of thy temple,'' and turns the psalm's thanks into a
request, \latin{Suscipiámus}, ``may we receive.'' It asks this so that we may
go before \latin{reparatiónis nostræ ventúra solémnia}, ``the coming
solemnities of our restoration,'' with fitting honours; Schuster takes them to
be Christmas (p.~323). Its noun \latin{misericórdia} is the Alleluia's: there
the Church asks God to show his mercy, and here she asks to receive it.
\zlabel{triptych:concise:themes:end}
